\section{Conditioned trees and uniform path sums}
\label{tree:section}

We isolate the probabilistic statement needed to pass from the block tree
to graph distances. Conditioned-tree height limits originate in
\cite{Aldous1993}; their stable-domain extension is developed in
\cite{Duquesne2003}. Joint depth-first process convergence to the same
excursion was proved by Marckert and Mokkadem under stronger moment
conditions~\cite{MarckertMokkadem2003}. We invoke the finite-variance,
attainable-size formulation of Broutin and
Marckert~\cite{BroutinMarckert2014} for the tree geometry.  We then prove a
uniform law of large numbers for marked ancestral paths, using only a
finite second offspring moment.

Let $\xi$ have distribution $(p_k)_{k\geq0}$ supported on the even
nonnegative integers, with
\begin{equation}
 p_0p_2>0,\qquad \E\xi=1,\qquad
 0<\sigma^2:=\operatorname{Var}(\xi)<\infty.
 \label{tree:assumptions}
\end{equation}
Write $\mathcal T$ for the corresponding plane Galton--Watson tree and
$\mathcal T_N$ for its law conditional on having $N$ nodes.  Its attainable
sizes are exactly the odd integers: every tree satisfies
$N-1=\sum_v\deg^+(v)$, and every odd size can be realized using only
offspring numbers zero and two.  All limits in this section run through
these odd integers.  Given the offspring numbers, attach independent
marks $A_v$ to the nodes, with a fixed law $\mathcal L_k$ at a node with
$k$ children.

Index the nodes in preorder by $v_0,\ldots,v_{N-1}$.  Let $H(v)$ denote
depth, put $H_j=H(v_j)$ for $j<N$ and $H_N=0$, and define the
exploration walk by
\begin{equation}
 W_0=0,\qquad W_{j+1}-W_j=\deg^+(v_j)-1
 \quad(0\leq j<N).
 \label{tree:walk-definition}
\end{equation}
Thus $W_j\geq0$ for $j<N$ and $W_N=-1$.  Interpolate both sequences
linearly between integer times, and write $W(v_j)=W_j$.

\begin{proposition}[Joint coding limit]
\label{tree:excursion}
Under \eqref{tree:assumptions},
\begin{equation}
 \left(\frac{W_{Nt}}{\sigma\sqrt N},
       \frac{\sigma H_{Nt}}{2\sqrt N}\right)_{0\leq t\leq1}
 \ \Longrightarrow\ (e(t),e(t))_{0\leq t\leq1}
 \quad\hbox{in }C([0,1],\R^2),
 \label{tree:walk-limit}
\end{equation}
where $e$ is a standard normalized Brownian excursion.
\end{proposition}
\begin{proof}
Broutin and Marckert's Theorem~3~\cite{BroutinMarckert2014} gives the
joint height and walk limit for uniform plane trees whose degree
frequencies and second moments converge to a critical finite-variance
law, with maximum degree $o(\sqrt N)$.  Their Lemma~11 verifies these
hypotheses in probability for conditioned critical Galton--Watson trees;
their Section~6 explicitly takes all attainable sizes.

Given the offspring counts $(n_k)$, our tree is uniform, since each
shape has probability $\prod_kp_k^{n_k}$.  To apply the deterministic
theorem to these random counts, take any subsequence and then a further
subsequence on which its degree hypotheses hold almost surely.
Conditional expectations of bounded continuous tests converge by the
theorem; dominated convergence removes the conditioning.  This proves
the limit along the full admissible sequence.

The source parametrizes height by $(N-1)t$.  Continuous-path tightness
makes the change to $Nt$ negligible before the last interval.  On that
interval, $H_{N-1}/\sqrt N\to0$ in probability because the limiting
excursion ends at zero, so appending $H_N=0$ is also harmless.  The marks
do not affect the tree-shape law.
\end{proof}

Put $\nu=\sigma^2/2$ and let $\Delta_N$ be the largest offspring number.

\begin{corollary}
\label{tree:maxdegree}
We have
\begin{equation}
 \max_v H(v)=O_{\P}(\sqrt N),\qquad
 \max_v W(v)=O_{\P}(\sqrt N),\qquad
 \Delta_N=o_{\P}(\sqrt N),
 \label{tree:degree-bound}
\end{equation}
and
\begin{equation}
 \max_v|W(v)-\nu H(v)|=o_{\P}(\sqrt N).
 \label{tree:walk-height-estimate}
\end{equation}
\end{corollary}
\begin{proof}
The height and walk bounds follow from Proposition~\ref{tree:excursion}.
Since both normalized processes converge jointly to the same excursion,
their uniform difference tends to zero in probability, proving
\eqref{tree:walk-height-estimate}.  Continuous-path tightness also makes
the largest walk increment divided by $\sqrt N$ tend to zero.  An
offspring number is a walk increment plus one.
\end{proof}

\begin{lemma}[Cost of conditioning]
\label{tree:conditioning}
As $N$ tends to infinity through odd integers,
\begin{equation}
 \pi_N:=\P(|\mathcal T|=N)
 \sim\frac{2}{\sigma\sqrt{2\pi}}N^{-3/2}.
 \label{tree:size-probability}
\end{equation}
\end{lemma}
\begin{proof}
Let $X_1,X_2,\ldots$ be independent with law $\xi-1$, and put
$S_r=\sum_{j=1}^rX_j$.  The exploration fills one waiting slot and
creates $\xi$ new ones at each step.  Thus total size $N$ means that
$S$ first reaches $-1$ at $N$.  A list of increments at least $-1$
totaling $-1$ has exactly one circular shift with this property: extend
the list periodically, decreasing its height by one per period, and
observe that its strict descending record endpoints occur once per
period.  Cyclic symmetry therefore gives the Otter--Dwass formula
(see also \cite[Theorem~15.5]{Janson2012})
\[
 \pi_N=\frac1N\P(S_N=-1).
\]
The lattice span of $X$ is exactly two because $-1$ and $1$ both have
positive probability.  At compatible sites $j\equiv r\pmod2$, the
finite-variance local limit estimate is
\[
 \P(S_r=j)=\frac{2}{\sigma\sqrt{2\pi r}}
             \exp\left(-\frac{j^2}{2r\sigma^2}\right)
             +o(r^{-1/2}),
\]
uniformly in $j$.  One can obtain this directly by Fourier inversion on
$[-\pi/2,\pi/2]$ with prefactor $1/\pi$: the characteristic function is
$1-\sigma^2\theta^2/2+o(\theta^2)$ near zero and has modulus strictly
less than one elsewhere on that interval.  Rescaling $\theta$ by
$r^{-1/2}$ gives the Gaussian integral, with uniformly vanishing
integrated error.  Taking $r=N$ and $j=-1$ proves the assertion.
\end{proof}

\subsection{A spine law and bounded labels}

For each strict ancestor of a node $v$, let $k$ be the ancestor's number
of children and $i\in\{1,\ldots,k\}$ the child leading toward $v$.
The waiting slots in the preorder exploration give
\begin{equation}
 W(v)=\sum_{\operatorname{path}(v)}(k-i),
 \qquad
 \widehat W(v):=\sum_{\operatorname{path}(v)}(i-1).
 \label{tree:waiting-slots}
\end{equation}
The second quantity is the preorder walk value at the corresponding
node after reversing every child order.  Reflection preserves the law
of $\mathcal T_N$ and preserves depths.  Hence
Corollary~\ref{tree:maxdegree} also gives
\begin{equation}
 \max_v|\widehat W(v)-\nu H(v)|=o_{\P}(\sqrt N).
 \label{tree:mirror-estimate}
\end{equation}
Only the tree shape is used here; no symmetry of the marks is required.

The size-biased spine construction is classical; see
Lyons, Pemantle, and Peres~\cite[Section~2]{LyonsPemantlePeres1995}.
For the marks and chosen child used here, define a probability law on
triples $(K,I,A)$ by
\begin{equation}
 \P_{\mathrm{sp}}(K=k,I=i,A\in B)=p_k\mathcal L_k(B),
 \qquad 1\leq i\leq k.
 \label{tree:spine-law}
\end{equation}
The total mass is $\E\xi=1$, and
\begin{equation}
 \E_{\mathrm{sp}}(K-I)=\frac{\E[\xi(\xi-1)]}{2}=\nu,
 \qquad \E_{\mathrm{sp}}K=2\nu+1<\infty.
 \label{tree:spine-means}
\end{equation}
For independent triples $Z_1,\ldots,Z_h$ with this law and every
nonnegative path function $F$, independence in the unconditioned tree
gives the path-counting identity
\begin{equation}
 \E\left[\sum_{H(v)=h}F(\operatorname{path}(v))\right]
 =\E_{\mathrm{sp}}F(Z_1,\ldots,Z_h).
 \label{tree:spine-identity}
\end{equation}
Indeed, at each ancestor one sums over its possible child choices; the
probability of each pair $(k,i)$, including its independent mark, is
exactly the factor in \eqref{tree:spine-law}.

For a label $b=b(k,i,A)$, set
\[
 B_b(v)=\sum_{\operatorname{path}(v)}b(k,i,A),\qquad
 \mu_b=\E_{\mathrm{sp}}b(K,I,A)
\]
whenever the expectation exists.

\begin{lemma}[Bounded path sums]
\label{tree:bounded-paths}
For every fixed bounded measurable real label $b$,
\begin{equation}
 \max_{v\in\mathcal T_N}|B_b(v)-\mu_bH(v)|
 =o_{\P}(\sqrt N).
 \label{tree:bounded-sums}
\end{equation}
\end{lemma}
\begin{proof}
Fix $C,\varepsilon>0$.  The exponential estimate for bounded independent
variables gives, uniformly for $h\leq C\sqrt N$,
\[
 \P_{\mathrm{sp}}\left(
 \left|\sum_{r=1}^h b(Z_r)-\mu_bh\right|>\varepsilon\sqrt N
 \right)\leq C' e^{-c'\sqrt N}.
\]
For short paths the event may be empty; otherwise the usual bound
$2\exp(-c\varepsilon^2N/h)$ gives the displayed estimate.  The
probability under $\mathcal T_N$ that any such path violates the desired
bound is at most
\[
 \pi_N^{-1}\sum_{h\leq C\sqrt N}C'e^{-c'\sqrt N}=o(1),
\]
by the unconditioned identity \eqref{tree:spine-identity} and
Lemma~\ref{tree:conditioning}.  Now let $C$ increase, using maximum-depth
tightness from Corollary~\ref{tree:maxdegree}.
\end{proof}

\subsection{Labels bounded by the offspring number}

We now allow labels that grow with the offspring number. The two
exploration orders control the sum of offspring numbers along every
ancestral path. Subtracting its bounded truncation will control the
remaining tail uniformly.

\begin{theorem}[Uniform marked path sums]
\label{tree:paths}
Let $b=b(k,i,A)$ be a fixed measurable label with $0\leq b\leq k$.
Then $\mu_b<\infty$ and
\begin{equation}
 \max_{v\in\mathcal T_N}|B_b(v)-\mu_bH(v)|=o_{\P}(\sqrt N),
 \qquad
 \max_{v\in\mathcal T_N}
 \left|B_b(v)-\frac{\mu_b}{\nu}W(v)\right|=o_{\P}(\sqrt N).
 \label{tree:marked-sums}
\end{equation}
\end{theorem}
\begin{proof}
The exact identities \eqref{tree:waiting-slots} give
\begin{equation}
 \sum_{\operatorname{path}(v)}k=W(v)+\widehat W(v)+H(v).
 \label{tree:total-label}
\end{equation}
By \eqref{tree:walk-height-estimate} and \eqref{tree:mirror-estimate},
this is $(2\nu+1)H(v)+o_{\P}(\sqrt N)$ uniformly in $v$.  For each
fixed $R$, Lemma~\ref{tree:bounded-paths} applies to the bounded label
$k\1_{\{k\leq R\}}$.  Subtracting its estimate yields
\begin{equation}
 \sum_{\operatorname{path}(v)}k\1_{\{k>R\}}
 =m_RH(v)+o_{\P}(\sqrt N)
 \quad\hbox{uniformly in }v,
 \qquad m_R:=\E_{\mathrm{sp}}[K\1_{\{K>R\}}].
 \label{tree:tail-estimate}
\end{equation}
The error statement here is for each fixed $R$.  Integrability of $K$
gives $m_R\to0$, and maximum-depth tightness consequently implies
\begin{equation}
 \lim_{R\to\infty}\limsup_{N\to\infty}
 \P\left(\max_v\sum_{\operatorname{path}(v)}
 k\1_{\{k>R\}}>\varepsilon\sqrt N\right)=0
 \quad(\varepsilon>0).
 \label{tree:uniform-tail}
\end{equation}
Apply Lemma~\ref{tree:bounded-paths} to $b\1_{\{k\leq R\}}$.
The discarded path sum is bounded by the tail in
\eqref{tree:uniform-tail}, and the discarded spine mean is bounded by
$m_R$.  Letting first $N$ and then $R$ tend to infinity proves the first
assertion.  The second follows from
\eqref{tree:walk-height-estimate} and $\nu>0$.
\end{proof}

This truncation uses $\E_{\mathrm{sp}}K=\E\xi^2<\infty$; it does not
require a second moment of the size-biased variable $K$, and hence does
not impose a third offspring moment.

We finish with a deterministic comparison for later use.

\begin{lemma}[Common ancestors and walk minima]
\label{tree:lca}
If $j\leq j'$ and $v_*$ is the lowest common ancestor of $v_j,v_{j'}$,
then
\begin{equation}
 W(v_*)\leq\min_{j\leq r\leq j'}W_r\leq W(v_*)+\Delta_N.
 \label{tree:lca-walk}
\end{equation}
\end{lemma}
\begin{proof}
Every intervening node is a descendant of $v_*$.  Its walk value is at
least $W(v_*)$ by \eqref{tree:waiting-slots}.  If $v_*=v_j$, the
upper bound is immediate.  Otherwise the interval contains the child of
$v_*$ leading toward $v_{j'}$; its walk value exceeds $W(v_*)$ by at
most the offspring number of $v_*$.
\end{proof}
